\documentclass[10pt,a4paper]{amsart}
\usepackage[margin=19mm]{geometry}
\usepackage{amsmath,amssymb,mathtools}
\usepackage[scr=rsfs,cal=euler]{mathalfa}
\usepackage{xfrac}
\usepackage[unicode,hidelinks,bookmarks=false]{hyperref}
\newcommand{\ie}{i.e.\ }
\newcommand{\cf}{cf.\ }
\newcommand{\resp}{respectively\ }
\newcommand{\Subsection}{\subsection}
\providecommand{\nobreakdash}{\nobreak\hbox{-}}
\setlength{\emergencystretch}{2em}
\multlinegap0pt
\usepackage{algorithm}\usepackage{algpseudocode}
\usepackage{mathtools}
\usepackage{amssymb}
\usepackage{float}
\usepackage{xcolor}
\newtheorem{lemma}{Lemma}
\newcommand{\C}[1]{\mathcal C(#1)}
\newcommand{\NN}{\mathbb N}
\newcommand{\bits}{\mathrel{\triangleleft}}
\title{A 4AP-free permutation of the positive integers}
\author{Boon Suan Ho}
\date{Source: arXiv:2609.12780v1 (CC BY 4.0)}
\begin{document}
\maketitle

\noindent\emph{Source and licence.} A 4AP-free permutation of the positive integers. Version arXiv:2609.12780v1 (CC BY 4.0), distributed by arXiv under the Creative Commons Attribution 4.0 International licence, http://creativecommons.org/licenses/by/4.0/, as recorded in the arXiv licence metadata for this exact version and shown on the abstract page https://arxiv.org/abs/2609.12780v1. Full licence notice: \texttt{licenses/arxiv-2609.12780-CC-BY-4.0.txt}.\par\smallskip

\noindent\textit{Editorial scope.} Counted unit: the article's lemma lem:extend (source0.tex lines 163--166). The article's complete original proof of it is delivered with it (source0.tex lines 168--314, one \texttt{proof} environment of 147 lines). No mathematics is written, repaired or re-derived: the statement and the proof are the article's own lines, in the article's own notation, and no retained line is substituted or re-flowed.  Retained uncounted as the source-local context the proof needs: lines 135--137.  Nothing was omitted from any retained span, so the package carries no omission record.  No retained line contains a citation, so the wrapper carries no bibliography and no bibliography entry.  No retained line contains a figure environment, an image-inclusion command or drawing code, so no image is delivered and the package declares no asset dependency.  Editorial formatting only, in addition: an AMS \texttt{amsart} preamble that restates the article's own declarations and macros used by the retained lines, namely the environments \texttt{lemma} on separate counters, together with the standard packages those lines need.  The retained environments print in this document as Lemma 1 in order of appearance, whereas the article prints the same items with the numbering of its own full document.  The complete native source of the article as distributed in the arXiv e-print for this exact version is bundled unchanged as \texttt{sources/210151/source0.tex}.
\begin{lemma}\label{lem:base}
Every finite set, listed in reverse $\bits$-order, is a safe word.
\end{lemma}

\begin{lemma}[Extension]\label{lem:extend}
Every safe finite word $P$ is an initial segment of a safe finite word
containing any prescribed finite set $T\subset\NN_0$.
\end{lemma}

\begin{proof}
We induct on the largest entry $m$ of $P$, taking $m=0$ if $P$ is empty. For
each value of~$m$, we require the statement to hold for every finite target set
$T$.

If $m=0$, then $P$ is either empty or consists only of $0$. List the elements of
$P\cup T$ without repetition in reverse $\bits$-order. Since $0$ is the
greatest element of $\bits$, this word begins with $P$, and it is safe by
Lemma~\ref{lem:base}.

Now suppose $m>0$. Split $P$ according to parity. Let $P_0$ be the word obtained
from the even entries of $P$ by dividing them by $2$, and let $P_1$ be obtained
from the odd entries by applying $2x+1\mapsto x$. Both $P_0$ and $P_1$ are safe.
Indeed, restricting $\C P$ to the even integers leaves first the even entries
of $P$, followed by the unused even integers in $\bits$-order. After applying
$2x\mapsto x$, the first block becomes $P_0$, while the second becomes
$\NN_0\setminus P_0$ in $\bits$-order, by the self-similarity of $\bits$.
Thus the resulting order is exactly $\C{P_0}$.
Similarly, after restricting $\C P$ to the odd integers and rescaling by
$2x+1\mapsto x$, we obtain exactly $\C{P_1}$. Since a restriction of a 4AP-free
order is 4AP-free, both $P_0$ and $P_1$ are safe. Moreover, every entry of
either word is at most $\lfloor m/2\rfloor<m$, so the induction hypothesis
applies to both.

\newpage
\begin{algorithm}[H]
\small
\hrule height .8pt
\kern 3pt
\noindent$\textsc{Extend}(P,T)$\par
\kern 3pt
\hrule
\begin{algorithmic}[1]
\Require A safe finite word $P$ and a finite target set $T\subset\NN_0$.
\Ensure A safe word $Q$ having $P$ as a prefix and containing every element of $T$.
\State $m\gets\max P$, with $m=0$ if $P$ is empty
\If{$m=0$}
  \State \Return the elements of $P\cup T$ in reverse $\bits$-order
\EndIf
\Statex \textit{Normalize the even and odd subsequences of $P$.}
\State $P_0\gets$ the even entries of $P$, in their original order, divided by $2$
\State $P_1\gets$ the odd entries of $P$, in their original order, mapped by $2x+1\mapsto x$
\Statex \textit{Extend the even part to cover the even targets.}
\State $T_0\gets\{x\in\NN_0:2x\in T\}$
\State $R_0\gets\Call{Extend}{P_0,T_0}$
\State $E_0\gets$ the suffix of $R_0$ obtained by deleting its prefix $P_0$
\State $E\gets E_0$ rescaled by $x\mapsto2x$
\State $h\gets\max E$, with $h=0$ if $E$ is empty
\Statex \textit{Extend the odd part, forcing enough small odd numbers to precede $E$.}
\State $T_1\gets\{x\in\NN_0:2x+1\in T\}\cup\{x\in\NN_0:x<h\}$
\State $R_1\gets\Call{Extend}{P_1,T_1}$
\State $O_1\gets$ the suffix of $R_1$ obtained by deleting its prefix $P_1$
\State $O\gets O_1$ rescaled by $x\mapsto2x+1$
\Statex \textit{Splice the new odd entries before the new even entries.}
\State \Return $P\,O\,E$
\end{algorithmic}
\hrule
\caption[Safe extension procedure $\textsc{Extend}(P,T)$]{%
The recursive construction used in the proof. The two recursive calls are on
words whose largest entries are strictly smaller than $m$. The extra targets
$x<h$ in the odd call force every odd integer at most $2h$ which is not already
in $P$ to be placed before the new even block $E$; doing this prevents
mixed-parity 4APs from forming.}\label{alg:extend}
\end{algorithm}

First apply the induction hypothesis to $P_0$, asking that it contain all
numbers $t/2$ with $t\in T$ even. After rescaling by $x\mapsto2x$, the result
consists of the old even subsequence of $P$, followed by some new even entries;
let $E$ denote the suffix consisting of these new entries. Let $h$ be the
largest entry of $E$, or set $h=0$ if $E$ is empty.

We now deal with the odd entries. Apply the induction hypothesis to $P_1$,
asking that it contain both the normalized odd elements $(t-1)/2$ of $T$ and all
nonnegative integers less than $h$. After rescaling by $x\mapsto2x+1$, the
result consists of the old odd subsequence of $P$, followed by a word $O$ of new
odd entries. In particular, every odd integer at most $2h$ either already occurs
in $P$ or occurs in $O$.

We claim that
$$
Q=P\,O\,E
$$
is the desired extension. By construction, the entries of $O$ and $E$ are new,
and they have opposite parity, so $Q$ has no repeated entries. It begins with
$P$ and contains every element of~$T$. It remains only to prove that $Q$ is
safe.

By construction, after restricting $\C Q$ to the even integers and rescaling
by $2x\mapsto x$, we obtain the safe extension of $P_0$ constructed above,
followed by all remaining nonnegative integers in $\bits$-order. Since that
extension of $P_0$ is safe, this order is 4AP-free. Likewise, after restricting
$\C Q$ to the odd integers and rescaling by $2x+1\mapsto x$, we obtain the safe
extension of $P_1$ constructed above, followed by all remaining nonnegative
integers in $\bits$-order, and hence this order is also 4AP-free. Thus both
parity restrictions of $\C Q$ are 4AP-free. Furthermore, after the old prefix
$P$, the order $\C Q$ has the form
$$
O,\quad E,\quad \text{unused odd integers},\quad
\text{unused even integers},
$$
because all odd integers precede all even integers under $\bits$.

It follows that
\begin{equation}\label{eq:odd-even}
\begin{matrix}
\textsl{if $x$ and $y$ do not occur in $P$, $x$ is even, $y$ is odd,
and $y\le2x$,}\\\textsl{then $y$ precedes $x$ in $\C Q$.}
\end{matrix}
\end{equation}
Indeed, if $x\in E$, then $x\le h$, so every odd $y\le2x$ which does not already
occur in $P$ belongs to $O$. If $x$ is still unused, then every odd integer
occurs before it.

Suppose for contradiction that a 4AP
$$
a,b,c,d
$$
occurs in this order in $\C Q$. If its common difference is even, then all four
terms have the same parity, contradicting the 4AP-freeness of the corresponding
parity restriction. Hence the common difference is odd, so the parities of
$a,b,c,d$ alternate.

First observe that $c$ does not occur in $P$. Otherwise $a,b,c$ all occur in
$P$, since $P$ is an initial segment of $\C Q$. The term $d$ would also follow
$c$ in the old completion $\C P$, whether or not $d$ occurs in $P$. Thus
$a,b,c,d$ would already form a 4AP in $\C P$, contradicting the safety of $P$.
Hence neither $c$ nor $d$ occurs in $P$.

If $c$ is even, then $d$ is odd. Since $a,b,c,d$ is an arithmetic progression,
$$
d=2c-b\le2c.
$$
By~\eqref{eq:odd-even}, $d$ must precede $c$, a contradiction.

It remains to consider the case where $c$ is odd, so that $b$ is even. We claim
that $b$ does not occur in $P$. Suppose otherwise. Since $P$ is an initial
segment of $\C Q$ and $a$ precedes $b$, we then also have $a\in P$. Since we
already know that $c,d\notin P$, in the old completion $\C P$ the terms $a,b$
occur in the initial block $P$, while $c,d$ occur in the $\bits$-ordered tail.
Since $c$ is odd and $d$ is even, $c$ precedes $d$ in that tail. Hence $a,b,c,d$
occur in this order in $\C P$, contradicting the safety of $P$.

Thus $b,c$ do not occur in $P$. But $c=2b-a\le2b$, so~\eqref{eq:odd-even} says
that the odd number $c$ precedes the even number $b$ in $\C Q$, contradicting
the assumed order of the progression in which $b$ precedes $c$.
Therefore $\C Q$ is 4AP-free, and $Q$ is safe. This completes the induction.
\end{proof}

\end{document}
